\section{The Bible of Hellenistic Jews and of the Apostles}

For perhaps three centuries the Greek Bible was simply the Scripture of
Greek-speaking Judaism. Philo, who shows no sign of reading Hebrew, built
an entire exegetical corpus on the Greek text and, as seen above, ranked it
with the original; the Pharos festival he describes was a Jewish feast, not
a Christian one. The deep vocabulary decisions of the third-century
translators---\textit{kyrios} for the divine name, \textit{christos} for
``anointed,'' \textit{diath\=ek\=e} (``testament'') for covenant,
\textit{nomos} for Torah---had by the first century become the religious
Greek of the eastern Mediterranean.\endnote{Class~S synthesis over the
translation's standard vocabulary; the individual equivalences are
uncontested and visible in any concordance of the Greek Bible.}

The New Testament was written inside that language. Its authors, writing
Greek for communities whose Scripture was Greek, quote the Old Testament
most often---not always---in its Greek form, and sometimes their argument
depends on wording the Hebrew does not have. Hebrews 10:5 reads Psalm~40
(Greek Psalm~39) with the Greek's ``a body hast thou prepared me,'' where
the Hebrew text speaks instead of opened ears; the letter's argument about
the Incarnation rides on the Greek noun. Acts 15:16--17 has James carry the
council of Jerusalem with Amos 9:11--12 in its Greek form---``the remnant
of men,'' and all the nations who bear God's name, will seek the
Lord---where the Hebrew has Israel possessing the remnant of Edom.
Matthew 1:23 cites Isaiah 7:14 with the Greek
translator's \textit{parthenos}, ``virgin,'' the very rendering over which
Jew and Christian would part ways in the second
century.\endnote{Greek wording checked against Brenton's translation of
the received Greek text (Psalm 39:7; Amos 9:11--12) at ebible.org; the
New Testament citations at their stated loci. The general statement is
bounded: citation practice varies by author and passage, quotations were
often adapted or made from memory, and no single ``apostolic Bible'' with
fixed contents can be reconstructed from the quotations. No percentage
of ``Septuagint quotations'' is asserted; the evidence does not support
one.}

Two inferences are commonly drawn from this and must be kept apart. That
the apostolic churches received the Jewish Greek Bible as Scripture is a
fact of the texts. What that reception implies for the boundaries of the
canon is a separate and harder question---taken up below---because the
first-century Greek Bible was not a bound volume with a table of contents
but a body of scrolls whose edges no surviving Jewish or Christian source
of that century defines.\endnote{The judgment that no fixed ``Alexandrian
canon'' is attested for the first century is a scholarly commonplace but
still a judgment (class~K/S); the fourth-century codices, sometimes cited
as evidence for it, are Christian artifacts five centuries after
Philadelphus and are treated below.}
